\begin{proof}

\begin{enumerate}[label={$\langle1\rangle$\arabic*.},leftmargin=*]
\item The exponential sequence and the two integral vanishings give
\[
 \operatorname{Pic}(X)=\operatorname{Pic}^0(X)
 \simeq H^1(X,\mathcal O_X).
\tag{W.54}\hypertarget{tag-03HP}{}\label{wbII:eq:W:5.27}
\]
Thus this Picard group is a connected irreducible complex vector group.
Choose additive Cech cocycles $a^{(r)}_{ij}$ representing a basis of
$H^1(X,\mathcal O_X)$. For coordinates $\lambda_r$ on that vector space,
\[
 g_{ij}(\lambda)=\exp\!\left(\sum_r\lambda_r a^{(r)}_{ij}\right)
\]
satisfy the multiplicative cocycle condition and define a holomorphic
line bundle on $X\times H^1(X,\mathcal O_X)$. Its restriction at every
parameter represents the corresponding line bundle, so proper
semicontinuity below applies to an actual holomorphic family.  Any fibre divisor \(D\) of
\(f\) is nonzero and effective, and its line bundle lies in
\(\operatorname{Pic}^0(X)\) by \(H^2(X;\mathbb Z)=0\).  Applying
\cref{wbII:lem:cdp-generic-endpoints} to \(B=T_X\) shows that the degree-zero and
degree-three groups in \eqref{wbII:eq:W:5.25} vanish away from two proper analytic subsets
of \(\operatorname{Pic}(X)\).

\item The cohomology-jump locus
\[
 \{M\in\operatorname{Pic}(X):h^2(X,T_X\otimes M)\ge1\}
\]
is analytic by upper semicontinuity.  It is proper because \eqref{wbII:eq:W:5.24} gives a
point outside it.  Its complement therefore meets the two endpoint
vanishing loci in a dense open set.  This proves \eqref{wbII:eq:W:5.25}, and on that set
\[
 \chi(X,T_X\otimes M)=-h^1(X,T_X\otimes M)\le0.
\tag{W.55}\hypertarget{tag-03HQ}{}\label{wbII:eq:W:5.28}
\]

\item Every \(M\in\operatorname{Pic}(X)\) has rationally trivial first Chern
class.  The hypotheses also give \(c_1(T_X)=0\) in rational cohomology and,
by Poincar\'e duality, \(H^4(X;\mathbb Q)=0\), so
\(c_2(T_X)=0\) rationally.  The degree-six term of
\(\operatorname{ch}(T_X\otimes M)\operatorname{td}(X)\) is consequently
\(\tfrac12c_3(T_X)\).  Hirzebruch--Riemann--Roch turns \eqref{wbII:eq:W:5.28} into
\[
 \chi(X,T_X\otimes M)
 =\frac12\left\langle c_3(T_X),[X]\right\rangle\le0,
\]
which proves \eqref{wbII:eq:W:5.26}.

\end{enumerate}
\end{proof}
